\documentclass{article}
\usepackage{amsmath,amssymb}
\usepackage[UTF8]{ctex}
\begin{document}

\section*{随机新辅助化疗 MILP (17)}

让$\mathcal D$表示可用的药物集合，$\mathcal Q$表示癌症细胞类型集合，而$\mathcal S=\{0,\ldots,S\}$则表示时间索引。对于固定的步长$h$，网格满足$t(0)=0$和$t(S)=T=Sh$的条件。场景集合由$k\in\{1,\ldots,K\}$进行索引。模型是……
\begin{align}
\min\quad &\sum_{q\in\mathcal Q}P_{q,S}^{(1)} \tag{17a}\\
\text{s.t.}\quad
&C_{d,s+1}=C_{d,s}-h\xi_d C_{d,s}+U_{d,s}/\mathcal V,
&&\forall d\in\mathcal D,\ s\in\{0,\ldots,S-1\}, \tag{17b}\\
&C_{d,0}=0, &&\forall d\in\mathcal D, \tag{17c}\\
&P_{q,s+1}^{(k)}=P_{q,s}^{(k)}+h\left(\Lambda\left(\ln(N_{q,\infty}^{(k)})-P_{q,s}^{(k)}\right)
-\sum_{d\in\mathcal D}\eta_{d,q}\exp(-\rho_{d,q}t(s))E_{d,s}\right),
&&\forall q\in\mathcal Q,\ s\in\{0,\ldots,S-1\},\ k\in\{1,\ldots,K\}, \tag{17d}\\
&P_{q,0}^{(k)}=\ln(N_{q,0}^{(k)}),
&&\forall q\in\mathcal Q,\ k\in\{1,\ldots,K\}, \tag{17e}\\
&E_{d,s}\geq 0, &&\forall d\in\mathcal D,\ s\in\mathcal S, \tag{10a}\\
&E_{d,s}\geq C_{d,s}-\beta_{d,\mathrm{eff}}, &&\forall d\in\mathcal D,\ s\in\mathcal S, \tag{10b}\\
&E_{d,s}\leq\beta_{d,\mathrm{conc}}Z_{E,d,s}, &&\forall d\in\mathcal D,\ s\in\mathcal S, \tag{10c}\\
&E_{d,s}\leq C_{d,s}-\beta_{d,\mathrm{eff}}+\beta_{d,\mathrm{conc}}(1-Z_{E,d,s}),
&&\forall d\in\mathcal D,\ s\in\mathcal S, \tag{10d}\\
&Z_{E,d,s}\in\mathbb B, &&\forall d\in\mathcal D,\ s\in\mathcal S, \tag{10e}\\
&P_{q,S}^{(k)}\leq P_{\mathrm{surg}}+
\ln\left(\frac{N_{q,0}^{(k)}}{\sum_{q\in\mathcal Q}N_{q,0}^{(k)}}\right)
+P_{q,\infty}(1-Z_{\mathrm{surg}}^{(k)}),
&&\forall q\in\mathcal Q,\ k\in\{1,\ldots,K\}, \tag{17g}\\
&\sum_{k=1}^{K}\mu^{(k)}Z_{\mathrm{surg}}^{(k)}\geq 1-\epsilon, && \tag{17h}\\
&Z_{\mathrm{surg}}^{(k)}\in\mathbb B, &&\forall k\in\{1,\ldots,K\}, \tag{17i}\\
&(\mathbf U,\mathbf C,\mathbf N_{\mathrm{neu}},\mathbf N_{\mathrm{lym}})\in\widehat{\mathcal O}. && \tag{17j}
\end{align}

\paragraph{情景与参数定义。}
肿瘤的异质性可以用一个离散随机向量来表示。
\[
\pi\in\mathbb R_+^{|\mathcal Q|},\qquad
\pi\in\{\pi^{(1)},\ldots,\pi^{(K)}\},\qquad
\mu^{(k)}=\Pr\{\pi=\pi^{(k)}\},\quad k\in\{1,\ldots,K\},
\]
根据所述的最可能情景前提，
\[
\mu^{(1)}\geq\mu^{(k)},\qquad \forall k\in\{1,\ldots,K\}.
\]
源代码中定义了$P_{\mathrm{surg}}=\ln(N_{\mathrm{surg}})$和$P_{q,\infty}=\ln(N_{q,\infty})$。这里，$N_{\mathrm{surg}}$是临床确定的固定可操作尺寸，而$\epsilon$则是未达到目标的概率。上标$(k)$表示在场景$k$下的数值；处理变量$U,C,E$则在不同场景中共享。目标是在场景1中最小化最终的对数种群，而不是对多个场景的期望值进行求和。

固定效应-隔室体积满足$\mathcal V>0$的条件。变量$U_{d,s}$、$C_{d,s}$和$E_{d,s}$分别表示给药剂量、浓度以及有效浓度。对数种群状态的含义为$P_{q,s}^{(k)}=\ln(N_{q,s}^{(k)})$。参数$\xi_d,\Lambda,\eta_{d,q},\rho_{d,q}$分别表示药物消除速率、Gompertz形状参数、癌细胞部分杀伤效应以及时间有效性衰减参数。论文中所述为$\eta_{d,q}\geq0$。固定的输入种群$N_{q,0}^{(k)}$和$N_{q,\infty}^{(k)}$分别对应于情景初始状态和渐近细胞类型种群。

\paragraph{选定的运筹学构建方法。}
在 (17j) 中，选择论文中的离散近似方法 (13)，而不是另一种 McCormick 约束方法 (12)。将 $N_{w,0}$ 写作固定预处理后的白细胞计数，将 $N_{w,s}$ 写作其离散化后的计数，将 $B_{d,s}$ 写作白细胞计数与延迟药物浓度乘积的近似值。论文假设……
\[
\beta_w=\min\left\{\frac{\beta_{\mathrm{neu}}}{\theta_{\mathrm{neu}}},
\frac{\beta_{\mathrm{lym}}}{\theta_{\mathrm{lym}}}\right\},\qquad
0\leq C_{d,s}\leq\beta_{d,\mathrm{conc}},\quad
\beta_w\leq N_{w,s}\leq N_{w,0},
\]
对于所有适用的$d\in\mathcal D$和$s\in\mathcal S$。延迟指数$\tau$对应于响应延迟$t_w$。在药物效应开始之前，第8页上未标注序号的重复内容
\[
N_{w,s+1}=N_{w,s}+h(\upsilon_w-\nu_wN_{w,s}),
\qquad \forall s\in\{0,\ldots,\tau-1\},
\]
其中，$N_{w,0}$表示初始计数。所选的延迟构造为
\begin{align}
N_{w,s+1}&=N_{w,s}+h\left(\upsilon_w-\nu_wN_{w,s}
-\sum_{d\in\mathcal D}\eta_{d,w}B_{d,s}\right),
&&\forall s\in\{\tau,\ldots,S-1\}, \tag{13a}\\
-\frac{\Delta}{2}\leq N_{w,s}-\left(\beta_w+\sum_{k=1}^{K}k\Delta Z_{w,s,k}\right)\leq\frac{\Delta}{2},
&&\forall s\in\{\tau,\ldots,S-1\}, \tag{13b--c}\\
&\sum_{k=0}^{K}Z_{w,s,k}=1, &&\forall s\in\{\tau,\ldots,S-1\}, \tag{13d}\\
&B_{d,s}=\sum_{k=1}^{K}(\beta_w+k\Delta)V_{d,s,k},
&&\forall d\in\mathcal D,\ s\in\{\tau,\ldots,S-1\}, \tag{13e}\\
&V_{d,s,k}\leq\beta_{d,\mathrm{conc}}Z_{w,s,k},
&&\forall d\in\mathcal D,\ s\in\{\tau,\ldots,S-1\},\ k\in\{1,\ldots,K\}, \tag{13f}\\
&V_{d,s,k}\leq C_{d,s-\tau},
&&\forall d\in\mathcal D,\ s\in\{\tau,\ldots,S-1\},\ k\in\{1,\ldots,K\}, \tag{13g}\\
&V_{d,s,k}\geq C_{d,s-\tau}+\beta_{d,\mathrm{conc}}(Z_{w,s,k}-1),
&&\forall d\in\mathcal D,\ s\in\{\tau,\ldots,S-1\},\ k\in\{1,\ldots,K\}, \tag{13h}\\
&V_{d,s,k}\geq0,
&&\forall d\in\mathcal D,\ s\in\{\tau,\ldots,S-1\},\ k\in\{1,\ldots,K\}, \tag{13i}\\
&Z_{w,s,k}\in\mathbb B,
&&\forall s\in\{\tau,\ldots,S-1\},\ k\in\{1,\ldots,K\}. \tag{13j}
\end{align}
在 (13b–c) 中的两个不等式，就是那对被打印出来的不等式。
\[
N_{w,s}-\left(\beta_w+\sum_{k=1}^{K}k\Delta Z_{w,s,k}\right)\leq\frac{\Delta}{2},\qquad
-N_{w,s}+\left(\beta_w+\sum_{k=1}^{K}k\Delta Z_{w,s,k}\right)\leq\frac{\Delta}{2}.
\]
第7页中关于离散近似的引言中，将$k\in\{0,\ldots,K\}$标记为$Z_{w,s,k}\in\mathbb B$；对于所有$k=0,\ldots,K$，都应包含该领域内容。打印版中的【13j】给出了更窄的范围$k\in\{1,\ldots,K\}$，这与引言中的描述存在冲突；应保留打印版中的【13j】范围，并指出这一冲突。

操作集中包含的中性粒细胞减少症和淋巴细胞减少症的约束条件包括：
\begin{align*}
N_{\mathrm{neu},s}&\geq\beta_{\mathrm{neu}},\qquad
N_{\mathrm{neu},s}=\theta_{\mathrm{neu}}N_{w,s},
&&\forall s\in\{0,\ldots,S\},\\
N_{\mathrm{lym},s}&\geq\beta_{\mathrm{lym}},\qquad
N_{\mathrm{lym},s}=\theta_{\mathrm{lym}}N_{w,s},
&&\forall s\in\{0,\ldots,S\}.
\end{align*}
这里的$\upsilon_w$代表白细胞的生成量，$\nu_w$代表白细胞的更新率，$\eta_{d,w}$代表白细胞的死亡率，而$\theta_{\mathrm{neu}},\theta_{\mathrm{lym}}$则代表总白细胞中的固定比例。固定的临床阈值为$\beta_{\mathrm{neu}},\beta_{\mathrm{lym}}$。药物对白细胞的效应取决于其浓度，即使浓度低于$\beta_{d,\mathrm{eff}}$，药物效果仍然有效，且不会因时间推移而减弱。

\paragraph{离散的运营管理限制。}
操作限制的定义将治疗期划分为$M$天、$D_m$天，以及$\mathrm{DAYS}=\{1,\ldots,M\}$天。而用餐时间则被标记为$\mathrm{MEALS}$。详细的官方附录A中，分别列出了卡培他滨、多西他赛和伊托比索的详细信息，这些内容被标记为$1,2,3$。对于卡培他滨而言：
\begin{align*}
C_{1,s}&\leq\beta_{1,\mathrm{conc}}, &&\forall s\in\mathcal S,\\
U_{1,s}&\leq\beta_{1,\mathrm{rate}}, &&\forall s\in\mathcal S,\\
\sum_{s\in D_m}U_{1,s}&\leq\beta_{1,\mathrm{cum}}, &&\forall m\in\mathrm{DAYS},\\
U_{1,s}&=0, &&\forall s\notin\mathrm{MEALS},\\
U_{1,s}&=\alpha_{1,\mathrm{pill}}Z_{1,\mathrm{pill},s}, &&\forall s\in\mathrm{MEALS},\\
Z_{1,\mathrm{pill},s}&\in\mathbb Z_+, &&\forall s\in\mathrm{MEALS}.
\end{align*}
对于多西他赛：
\begin{align*}
C_{2,s}&\leq\beta_{2,\mathrm{conc}}, &&\forall s\in\mathcal S,\\
U_{2,s}&\leq\beta_{2,\mathrm{rate}}, &&\forall s\in\mathcal S,\\
\sum_{s\in D_m}U_{2,s}&\leq\beta_{2,\mathrm{cum}}(1-Z_{2,\mathrm{rest}}^m), &&\forall m\in\mathrm{DAYS},\\
\sum_{l=0}^{\min\{\alpha_{2,\mathrm{rest}},M-m\}}(1-Z_{2,\mathrm{rest}}^{m+l})&\leq1, &&\forall m\in\mathrm{DAYS},\\
Z_{2,\mathrm{rest}}^m&\in\mathbb B, &&\forall m\in\mathrm{DAYS}.
\end{align*}
对于依托泊苷：
\begin{align*}
C_{3,s}&\leq\beta_{3,\mathrm{conc}}, &&\forall s\in\mathcal S,\\
U_{3,s}&\leq\beta_{3,\mathrm{rate}}, &&\forall s\in\mathcal S,\\
\sum_{s\in D_m}U_{3,s}&\leq\beta_{3,\mathrm{cum}}, &&\forall m\in\mathrm{DAYS},\\
U_{3,s}&=0, &&\forall s\notin\mathrm{MEALS},\\
U_{3,s}&=\alpha_{3,\mathrm{pill}}Z_{3,\mathrm{pill},s}, &&\forall s\in\mathrm{MEALS},\\
Z_{3,\mathrm{pill},s}&\in\mathbb Z_+, &&\forall s\in\mathcal S.
\end{align*}
源文档并未明确这些数学行中的具体用餐时间指标。多西他赛的休息终点值是按照打印格式包含的。休息相关的行适用于所述的需要休息的药物；口服药物的行则适用于所述的口服药物。

运算集还要求$U_{d,s},C_{d,s},E_{d,s},N_{w,s},N_{\mathrm{neu},s},N_{\mathrm{lym},s}$的非负性，不过癌症对数种群的状态$P_{q,s}^{(k)}$不必是非负的。所选的构造方式将$B_{d,s}$作为近似变量，而$V_{d,s,k}\geq0$则按照上述规定进行处理。为底层模型所设定的参数中包含了$\mathcal V>0$；论文中的欧拉稳定性结果还包含了$\xi_d,h,\mathcal V>0$和$h<2/\xi_d$，以及$\Lambda,h>0$和$h<2/\Lambda$。这些稳定性条件是这些分析结果的前提条件，而不是额外的行（17）。

\paragraph{印刷系数区分。}
方程(17d)打印的是与场景相关的$N_{q,\infty}^{(k)}$，而方程(17g)打印的是与场景无关的$P_{q,\infty}$，其定义为$\ln(N_{q,\infty})$。两种打印形式都被保留下来；方程(17g)中的系数并未被转换为与场景相关的量。概率约束替代方案(17g–(17i))是选定的保守构造方式；原始的概率约束并未作为额外的行被添加进来。

\end{document}